% Part: intuitionistic-logic
% Chapter: tableaux
% Section: introduction

\documentclass[../../../include/open-logic-section]{subfiles}

\begin{document}

\olfileid{int}{tab}{int}

\olsection{جان پہچان}

!!^{tableau}s !!{signed
  formula}s دے خاص درخت نیں، جیہناں دیاں ٹہنیاں تھلّے ول ودھدیاں نیں۔
یعنی ہر نوڈ اک جوڑی اے، سچائی دی قدر دا نشان ($\True$ یا $\False$) تے
!!a{sentence}:
\[
\sFmla{\True}{!A} \text{ یا } \sFmla{\False}{!A}.
\]
!!^a{tableau} کئی \emph{مفروضیاں} نال شروع ہوندا اے۔ ہر اگلا !!{signed formula}
استنباط دے قاعدیاں وچوں اک نوں لا کے بندا اے۔ کجھ قاعدے درخت دے کسے
سِرے اُتے اک یا ودھ !!{signed formula}s پا دیندے نیں۔ کجھ ہور دو نویں سِرے
بناندے نیں تے دو ٹہنیاں نکلدیاں نیں۔ قاعدے توں ملے !!{signed formula}s وچ
!!{formula} اوس !!{signed formula} دے فارمولے توں گھٹ گنجلدار ہوندا اے جیہدے
اُتے قاعدہ لگیا سی۔ جدوں ٹہنی وچ $\sFmla{\True}{!A}$ تے
$\sFmla{\False}{!A}$ دونیں ہون، اسیں ٹہنی نوں \emph{بند} آکھدے آں۔
جے !!a{tableau} دی ہر ٹہنی بند اے، تاں سارا !!{tableau} بند اے۔ بند
!!{tableau} اک !!a{derivation} اے جیہڑا وکھاندا اے کہ !!{tableau} دے شروع
وچ ورتے !!{signed formula}s دا سیٹ اکٹھا پورا ہون جوگا نہیں۔
ایس نال $\Proves$ تعلق دی تعریف کیتی جا سکدی اے۔
$\Gamma \Proves !A$ اُدوں تے صرف اُدوں اے جدوں اک محدود سیٹ~$\Gamma_0 = \{!B_1,
\dots, !B_n\} \subseteq \Gamma$ ہووے جیہدے لئی ہِیٹھلے مفروضیاں دا بند
!!{tableau} ہووے:
\[
\{\sFmla{\False}{!A}, \sFmla{\True}{!B_1}, \dots, \sFmla{\True}{!B_n}\}.
\]

وجدانیت پسند منطق لئی سانوں !!{signed formula} دی دھارنا ودھاؤنی پئے گی
تے رابطیاں دے قاعدے وی بدلنے پین گے۔ نشان ($\True$ یا $\False$) دے نال نال،
امکانی !!{tableau}s وچ !!{formula}s دے \emph{اگیتے}~$\sigma$ وی ہوندے نیں۔
اگیتے مثبت صحیح عدداں دیاں غیر خالی لڑیاں نیں، یعنی $\sigma \in
(\PosInt)^* \setminus \{\emptyseq\}$۔ اسیں ایہو جیہے اگیتے آلے دوالے
$\tuple{\ }$ توں بناں لکھدے آں، تے وکھرے وکھرے !!{element}s نوں~$.$ نال وکھ
کردے آں، $,$ نال نہیں۔ جے $\sigma$ اک اگیتا اے، تاں $\sigma.n$، $\sigma \concat \tuple{n}$
اے۔ مثلاً جے $\sigma = 1.2.1$، تاں $\sigma.3$، $1.2.1.3$ اے۔ سو مثال لئی،
\[
\sFmla{\True}{!A \lif (!B \lif !C)}[1.2]
\]
اک \emph{اگیتے والا !!{signed formula}} اے (یا مختصر طور اُتے
\emph{اگیتے والا !!{formula}})۔

سمجھن لئی اگیتے نوں نمونے وچ اک جہان دا ناں سمجھ لو۔ ایہ نمونہ
!!a{tableau} دی اک ٹہنی والے !!{formula}s پورے کر سکدا اے۔ جے $\sigma$
کسے جہان دا ناں اے، تاں $\sigma.n$ اوہ جہان دا ناں اے جیہڑا
~$\sigma$ ناں والے جہان توں رسائی وچ اے۔

وجدانیت پسند نمونیاں وچ رسائی دا تعلق انعکاسی تے متعدی اے۔ اگیتیاں دے
لحاظ نال ایہدا مطلب اے کہ $\sigma$، اپنے آپ~$\sigma$ توں رسائی وچ اے،
تے ہر اوہ اگیتا وی جیہڑا~$\sigma$ نوں ودھاوے، یعنی
$\sigma.n_1.\cdots.n_k$ شکل والا اگیتا۔ اسیں نشان $\sigma.*$ ورتاں گے،
جیہڑا خود $\sigma$ تے اوہدی ہر توسیع دسدا اے۔ دوجے لفظاں وچ
$\sigma.*$ اوہ سارے، تے صرف اوہی، اگیتے نیں جیہڑے~$\sigma$ توں رسائی وچ نیں۔

\end{document}
